\label{sec:process}
\subsection{The sticky localization}
\begin{definition}
Fix a centre $m$ and walls $\{w_a\cdot x=b_a\}$. Start from a resolved probe precision $A_0=a_0\mathrm{Id}$. Run Klartag's evolution on the precision,
\[
dA_t=\sigma\,\pi_t(dW_t),\qquad \pi_t=\text{projection onto }\{B:\ \langle y_a\otimes y_a,B\rangle=0\ \text{for every contact }a\},
\]
where $y_a=A^{-1}w_a$ is the direction to wall $a$'s tangency point (Proposition~\ref{prop:sliding}). A wall becomes a contact when it first
touches the probe, and stays one. The process freezes when the $y_a\otimes y_a$ span $\mathrm{Sym}_n$.
\end{definition}
By It\^o, $\E\log\det A_T=\log\det A_0-\frac{\sigma^2}2\E\int_0^T\delta_tdt$, with $\delta_t=\sum_{i,j}|\pi_t(u_i\otimes_su_j)|^2/(\lambda_i\lambda_j)$. Since $\log\det$ is
concave, the expected probe volume grows at a rate set by the free directions.
\emph{Check} ($n=3$, $14$ affine walls, $200$ paths, $T=2$):
\begin{itemize}[leftmargin=*]
\item $\E\log\det A_T=18.700\pm0.050$, against the It\^o prediction $18.683\pm0.010$, from $\log\det A_0=19.114$;
\item contacts accumulate to $1.9$ on average and up to $6=\dim\mathrm{Sym}_3$, i.e.\ perfect.
\end{itemize}
With homogeneous walls and the centre free, Proposition~\ref{prop:dilfree} keeps the dilation free forever. The probe then never
freezes, and its log-scale performs the free walk that is the critical mode.

\subsection{The maximal resolved probe and its code}
\begin{proposition}[Klartag's terminal state meets John; proved, and checked]\label{prop:john}
The locally maximal resolved probe about $m$ is the polar of the L\"owner ellipsoid of the dual wall points $\pm w_a/|w_a\cdot m-b_a|$. It touches
the contact walls at Mahalanobis distance $1$, and its precision satisfies the eutaxy (John) condition
\[
A=\sum_a\lambda_a\,x_a\otimes x_a,\qquad \lambda_a\ge0,\qquad \sum\lambda_a=n,\qquad x_a=\frac{w_a}{|w_a\cdot m-b_a|}.
\]
This is the polar form of Voronoi's eutaxy: for lattice points the inverse of the ellipsoid matrix is the frame sum, while for walls the
precision is. In the whitened frame $u_a=A^{-1/2}x_a$, $\sum_a\lambda_au_a\otimes u_a=\mathrm{Id}$: the contact normals form a weighted tight frame, a
spherical $2$-design.
\end{proposition}
\emph{Check} ($n=4$, $9$ walls): Khachiyan's algorithm gives $4$ active walls, all at Mahalanobis distance $1.0000$, with
$\|\sum\lambda_au_a\otimes u_a-\mathrm{Id}\|=1.3\times10^{-14}$ and $\sum\lambda_a=4.0000$.

This is the transformation that turns a local process into a code. Klartag's $T$ maps his terminal ellipsoid to the ball and his
lattice to a packing. Our whitening $A^{1/2}$ (the map $x\mapsto A^{1/2}(x-m)$) maps the maximal resolved probe at an input to the unit ball, and the walls around that
input to a tight frame of contact normals. The local geometry of linearity around an input, expressed in its own natural metric, is a
spherical design: the neighbourhood of every input carries a small frame code.

\subsection{Cells as cones: the conic Steiner formula and the global gate-count law}
\textsc{ReLU} is the metric projection onto the orthant, $\relu(z)=\Pi_{\R^n_+}(z)$. Moreau's decomposition $z=\Pi_K z+\Pi_{K^\circ}z$, i.e.\
$\relu(z)-\relu(-z)=z$, is the cone--polar duality, our analogue of lattice--dual lattice. For any closed convex cone $K$ and
$g\sim N(0,\mathrm{Id})$, the conic Steiner formula (McCoy--Tropp) says
\[
\E\,f\big(|\Pi_Kg|^2,|\Pi_{K^\circ}g|^2\big)=\sum_kv_k(K)\,\E f(\chi^2_k,\chi^2_{n-k}),
\]
where $v_k$ is the law of the dimension of the face that $\Pi_Kg$ lands in (the conic intrinsic volumes).
\begin{itemize}[leftmargin=*]
\item For the orthant, $v_k$ is binomial: the face is the gate pattern. For orthogonal weights the first layer is exactly this, so
$v_k$ is the global gate-count law.
\item \emph{Check}, in a layer-1 cell $C=W^{-1}\R^n_+$ with Gaussian $W$ ($n=10$, $2\times10^4$ samples, exact projections by NNLS):
\begin{itemize}
\item $\E|\Pi g|^2=3.179\pm0.020$ against $\sum kv_k=3.188$;
\item $\E|\Pi g|^4=17.89\pm0.24$ against $17.93$;
\item $\E|\Pi g|^2|\Pi^\circ g|^2=20.12\pm0.16$ against $20.33$.
\end{itemize}
\item The cell's statistical dimension is $3.19$, against $n/2=5$ for orthogonal weights. Gaussian weights make the cell thinner than an
orthant.
\end{itemize}

\subsection{Siegel's principle for the network's rows}
Klartag's analysis rests on Siegel's mean value theorem: a random lattice is, on average, Lebesgue measure. Our analogue is stage
12's self-localization. Through Gaussian rows independent of the state, the fields $r_a$ of a probe are, across rows, $N(0,t)$-distributed.
The ``theta series'' of the wall code relative to the probe therefore has the Siegel mean
\[
\Theta(s)=\sum_a e^{-sr_a^2/2},\qquad \E\,\Theta(s)=n\,(1+st)^{-1/2}.
\]
Its value at $s=1$ is $\sqrt{2\pi}$ times the wall density, $n/\sqrt{1+t}$, which is stage 12's wall-density law. Rogers' second-moment
formula has an analogue as well. Given the state, the rows are independent (Poisson-like), and the correlations come only through the
shared state. The leading one is the collective (dilation) mode, the Mattis component of stage 12, which is once more the charge.
